\documentclass[10pt,a4paper]{amsart}
\usepackage[margin=19mm]{geometry}
\usepackage{amsmath,amssymb,mathtools}
\usepackage[scr=rsfs,cal=euler]{mathalfa}
\usepackage{xfrac}
\usepackage[unicode,hidelinks,bookmarks=false]{hyperref}
\newcommand{\ie}{i.e.\ }
\newcommand{\cf}{cf.\ }
\newcommand{\resp}{respectively\ }
\newcommand{\Subsection}{\subsection}
\providecommand{\nobreakdash}{\nobreak\hbox{-}}
\setlength{\emergencystretch}{2em}
\multlinegap0pt
\usepackage{amssymb}
\usepackage{enumerate}
\usepackage{mathtools}
\usepackage{cancel}
\usepackage{algorithm}\usepackage{algpseudocode}
\newtheorem{theorem}{Theorem}
\newcommand{\Cc}{\mathcal{C}}
\newcommand{\Dc}{\mathcal{D}}
\newcommand{\Kc}{\mathcal{K}}
\DeclareMathOperator*{\argmin}{argmin}
\renewcommand{\bf}{\mathbf{f}} % NOTE: use \textbf if you really want to write bold non-math text.
\newcommand{\bg}{\mathbf{g}}
\newcommand{\bk}{\mathbf{k}}
\newcommand{\bu}{\mathbf{u}}
\newcommand{\bx}{\mathbf{x}}
\newcommand{\co}{{\operatorname{co}}}
\newcommand{\des}{{\operatorname{d}}}
\newcommand{\longthmtitle}[1]{\mbox{}{\textit{(#1):}}}
\newcommand{\map}[3]{#1:#2 \rightarrow #3}
\newcommand{\real}{\mathbb{R}}
\newcommand{\setdefb}[2]{\big\{#1 \; | \; #2\big\}}
\newcommand{\until}[1]{[#1]}
\title{\LARGE \textbf{Combinatorial Control Barrier Functions: Nested Boolean and $p$-choose-$r$ Compositions of Safety Constraints}}
\author{Pio Ong, Haejoon Lee, Tamas G. Molnar, Dimitra Panagou, Aaron D. Ames}
\date{Source: arXiv:2509.10716v2 (CC BY 4.0)}
\begin{document}
\maketitle

\noindent\emph{Source and licence.} \LARGE \textbf{Combinatorial Control Barrier Functions: Nested Boolean and $p$-choose-$r$ Compositions of Safety Constraints}. Version arXiv:2509.10716v2 (CC BY 4.0), distributed by arXiv under the Creative Commons Attribution 4.0 International licence, http://creativecommons.org/licenses/by/4.0/, as recorded in the arXiv licence metadata for this exact version and shown on the abstract page https://arxiv.org/abs/2509.10716v2. Full licence notice: \texttt{licenses/arxiv-2509.10716-CC-BY-4.0.txt}.\par\smallskip

\noindent\textit{Editorial scope.} Counted unit: the article's theorem thm:pCr\_safety (source0.tex lines 427--449). The article's complete original proof of it is delivered with it (source0.tex lines 450--471, one \texttt{proof} environment of 22 lines). No mathematics is written, repaired or re-derived: the statement and the proof are the article's own lines, in the article's own notation, and no retained line is substituted or re-flowed.  Retained uncounted as the source-local context the proof needs: lines 254--257; lines 261--263; lines 365--367; lines 418--421.  Nothing was omitted from any retained span, so the package carries no omission record.  The retained lines cite 3 works, whose entries are transcribed verbatim from the article's own bibliography source in the e-print and delivered with short editorial labels recorded in the item's reference mapping.  No retained line contains a figure environment, an image-inclusion command or drawing code, so no image is delivered and the package declares no asset dependency.  Editorial formatting only, in addition: an AMS \texttt{amsart} preamble that restates the article's own declarations and macros used by the retained lines, namely the environments \texttt{theorem} on separate counters, together with the standard packages those lines need.  The retained environments print in this document as Theorem 1 in order of appearance, whereas the article prints the same items with the numbering of its own full document.  The complete native source of the article as distributed in the arXiv e-print for this exact version is bundled unchanged as \texttt{sources/184992/source0.tex}.
\begin{equation}
    \label{sys:ctrl-affine}
    \dot\bx = \bf(\bx)+\bg(\bx)\bu
\end{equation}

\begin{equation}\label{eq:safeset}
    \Cc = \setdefb{\bx\in\real^n}{h(\bx)\geq 0}.    
\end{equation}

\begin{equation}\label{eq:sorting_CBFs}
    h(\bx)  = \max^r \{h_k(\bx)\}_{k=1}^p = \min^{p-r+1} \{h_k(\bx)\}_{k=1}^p.
\end{equation}

   \begin{equation}
    \dot h_k(\bx,\bu) > -\alpha \big( h(\bx)+\vert h_k(\bx)-h(\bx)\vert \big), \quad \forall k\in\until{p}.
    \label{eq:c-cbf}
    \end{equation}

\begin{theorem}\longthmtitle{Combinatorial Safety}
\label{thm:pCr_safety}
    Consider the system~\eqref{sys:ctrl-affine}.
    Let $h$ be constructed from a combination of primitive CBFs $\{h_k\}_{k=1}^p$ as in~\eqref{eq:sorting_CBFs}. If $h$ is a $p$-choose-$r$ CBF for~\eqref{sys:ctrl-affine}, then the set $\Cc$ in~\eqref{eq:safeset} is control invariant (safe).
    
    In addition, any state feedback controller ${\map{\bk}{\real^n}{\real^m}}$,  ${\bu=\bk(\bx)}$, that is continuous and satisfies:
    \begin{equation}
        \label{eq:pCr_BC}
    \dot h_k(\bx,\bk(\bx)) \geq -\alpha \big( h(\bx)+\vert h_k(\bx)-h(\bx)\vert \big), ~\forall k\in\until{p} 
    \end{equation}
    on a neighborhood ${\Dc\supset\Cc}$, renders the set $\Cc$ forward invariant and ensures the bounds:
    \begin{equation}\label{eq:pCr_bound}
        \frac{d}{dt} \big(\max^j\{h_k(\bx(t)\}_{k=1}^p\big) \geq -\alpha\big(\max^j\{h_k(\bx(t)\}_{k=1}^p\big)
    \end{equation}
    for all ${j \geq r}$  at almost every time $t\geq 0$. Consequently, for any initial condition~${\bx_0\in\Cc}$, there exist at least $r$ indices ${k\in \until{p}}$ at almost every time ${t\geq 0}$ such that ${h_k(\bx(t))\geq 0}$.

    In particular, the following CBF-QP:
    \begin{align}\label{eq:pCr-CBF-QP}
    \bk(\bx)  =\argmin_{\bu\in\real^m} \quad & \|\bu-\bk_\des(\bx)\|^2                                                                  \\
             \textup{s.t.} \quad  \dot h_k(\bx,\bu) &\geq -\alpha \big( h(\bx)+\vert h_k(\bx)-h(\bx)\vert \big), ~\forall k\in\until{p} \nonumber
    \end{align}
    is a continuous controller satisfying~\eqref{eq:pCr_BC}.
\end{theorem}

\begin{proof}
    Because each $\max^j\{h_k(\bx)\}_{k=1}^p$ is nonsmooth, the proof relies on nonsmooth barrier function theory. In particular, let $\Kc_j(\bx) = \setdefb{k\in\until p}{h_k(\bx) = \max^j\{h_k(\bx)\}_{k=1}^p}$ be the set of indices of CBFs with the same value as the function of interest.
    Then we have the subgradient~\cite{FHC:83}:
    \begin{align*}
    \partial \big(\max^j\{h_k(\bx)\}_{k=1}^p\big) &= \co\bigcup_{k\in\Kc_j(\bx)} \left\{\frac{\partial h_k}{\partial \bx}(\bx)\right\} \\
    \implies \frac{d}{dt} \big(\max^j\{h_k(\bx(t)\}_{k=1}^p\big)&\in \co\bigcup_{k\in\Kc_j(\bx)}\{\dot h_k(\bx(t),\bu(t))\}.
    \end{align*}
    The implication above follows from the nonsmooth chain rule~\cite[Thm. 2.3.10]{FHC:83} (see also, \cite{PG-JC-ME:17}).
    Then we note the inequalities in~\eqref{eq:pCr_BC} enforce for every $k\in\Kc_j(\bx)$:
    $$
    \dot h_k(\bx(t),\bu(t)) \geq -\alpha(h_k(\bx(t)))
    $$
    when $j\geq r$, from the definition of the sorting operator $\max^j$. Hence, the bounds given in~\eqref{eq:pCr_bound} follow. 
    
    Furthermore, when considering the case $j=r$, we derive:
    $$
    \langle\xi,~\bf(\bx)+\bg(\bx)\bk(\bx)\rangle \geq -\alpha(\max^r\{h_k(\bx)\}_{k=1}^p), 
    $$
    for all $\xi \in \partial\big(\max^r\{h_k(\bx)\}_{k=1}^p\big)$ on the neighborhood $\Dc$. The function $\max^r\{h_k(\bx)\}_{k=1}^p$ is therefore a nonsmooth barrier function~\cite[Prop. 2]{PG-JC-ME:17} for the closed-loop system, verifying the statement of forward invariance.
    
    To prove control invariance, we show that the CBF-QP is a valid safeguarding controller. The definition of the $p$-choose-$r$ CBFs using strict inequality in~\eqref{eq:c-cbf} ensures that the CBF-QP satisfies Slater's condition at every state $\bx$ in a neighborhood $\Dc\supset \Cc$ of the safe set, so it is well-defined and continuous there~\cite{PM-AA-JC:25}. In addition, the CBF-QP enforces~\eqref{eq:pCr_BC} from its constraint, so we may deduce safety.
\end{proof}

\begin{thebibliography}{99}
\bibitem[FHC:83]{FHC:83} F.~H. Clarke, {\em Optimization and Nonsmooth Analysis}. \newblock Canadian Mathematical Society Series of Monographs and Advanced Texts, New York: Wiley, 1983.
\bibitem[PG-JC-]{PG-JC-ME:17} P.~Glotfelter, J.~Cort\'es, and M.~Egerstedt, ``Nonsmooth barrier functions with applications to multi-robot systems,'' {\em IEEE Control Systems Letters}, vol.~1, no.~2, pp.~310--315, 2017.
\bibitem[PM-AA-]{PM-AA-JC:25} P.~Mestres, A.~Allibhoy, and J.~Cortés, ``Regularity properties of optimization-based controllers,'' {\em European Journal of Control}, vol.~81, p.~101098, 2025.
\end{thebibliography}
\end{document}
